\section{Deriving a weighted normal double from coorientation}
\label{sec:weighted-coorientation}

The first native topology-spectrum integration of
Section~\ref{sec:topology-spectra-native} used one boundary query and two
weighted surface queries. The ordinary component-count API already had a
finite sufficient certificate that the normal double is two copies of its
source (Section~\ref{sec:normal-coorientation}). This section extends that
certificate to the entire weighted histogram. A successful finite colouring
now removes the second weighted discovery and its replay, rather than merely
predicting its number of orbits. An unsuccessful colouring uses the complete
double query. The returned topology spectrum remains exact in both cases.

\subsection{A weight-preserving cover isomorphism}

The base normal interval system has points $j$ and the doubled system has
adjacent lifts $2j+s$, $s\in\{0,1\}$. A base arc pairing $\phi$ has reversal
bit $r$. Its doubled pairing is
\[
 2j+s\longmapsto2\phi(j)+(s\mathbin\oplus r).
\]
For a preserving map this follows by doubling its translation and interval
endpoints. For a reversing map, its doubled endpoint is $2d+1$ rather than
$2d$, and substitution reverses the sheet bit. The bit is still necessary
when the base interval contains just one point; forgetting it loses the
normal orientation cover.

Give each nonempty ambient-edge block a bit $\epsilon_u$. For every base
pairing from block $u$ to block $v$, require
\[
 \epsilon_u\mathbin\oplus\epsilon_v=r.
\]
These are the established block-constant coorientation constraints. Blocks
with no incident pairing may take bit zero. Failure is only failure of this
sufficient finite representation: an actual coorientation may vary within
one large block.

\begin{proposition}[Weighted double trivialization]
If these constraints are satisfied, let the base point weight be any vector
$w(j)\in\mathbb Z^d$ and suppose the doubled weight obeys
$w_2(2j+s)=w(j)$. If $H(v)$ is the base orbit histogram and $K(v)$ the
doubled histogram, then $K(v)=2H(v)$ for every weight vector $v$.
\end{proposition}
\begin{proof}
The corrected sheet label
$\sigma(j,s)=s\mathbin\oplus\epsilon_{\mathrm{block}(j)}$ is invariant
under every doubled pairing, by the displayed constraint. The map
$(j,s)\mapsto(j,\sigma(j,s))$ is a bijection from the doubled point set to
two labelled copies of the base set and conjugates each generator to the
base generator within a fixed label. Thus each doubled orbit maps
bijectively to one base orbit. By the point-weight hypothesis that bijection
preserves its sum. Each base orbit has one copy of each label, proving the
histogram identity, including zero or negative weight entries and binary
multiplicities.
\end{proof}

The multiplicity doubles; the signature does not. A base component of
signature $(\chi,b)$ gives two components of signature $(\chi,b)$,
not a component of signature $(2\chi,2b)$. The result is an interval-cover
theorem for arbitrary additive weights. The maintained application uses
$d=2$ and does not introduce a general new weighted-proof API.

\subsection{Both topology weights satisfy the hypothesis}

The Euler weight is the sum of point ownership for vertices, minus arc
ownership for surface edges, plus corner ownership for discs. Every original
half-open ownership interval $[a,b)$ becomes $[2a,2b)$ at scale two with
the same integer weight. Hence both lifts of every point receive its old
Euler contribution, even where many ownership intervals overlap.

For the second coordinate, the boundary transversal selects one original
point on each actual boundary circle. The embedding into the full point
universe preserves edge-block offsets. Its doubled interval selects both
adjacent lifts of that point. These are precisely the two lifted boundary
markers, so the boundary-marker point weight also pulls back unchanged.
No new doubled boundary transversal is required.

In an orientable ambient manifold the normal double is the orientation
cover. The colouring therefore also proves every source component
orientable. The usual identities $H=O+N$ and
$K(v)=2O(v)+N(v/2)$ then agree with $O=H$, $N=0$. The maintained sparse
inversion and its independent forward checks are still used, avoiding a
second special spectrum-restoration implementation. At the zero signature,
$N(0)=2H(0)-K(0)=0$ and $O(0)=H(0)$, so tori/Klein-bottle ambiguity is
resolved by the geometric certificate, not by an Euler-value heuristic.
The existing one-sided scaling and vertex-link restoration remain unchanged
on fallback inputs.

The finite colouring is never inferred from a primitive cohomology class,
from positive Euler characteristic, or from a primitive coordinate vector.
The disconnected primitive fibre of
Section~\ref{sec:completion-theorems} illustrates why source-specific
proofs matter. Even an orientable even multiple of a M\"obius band can fail
the block-constant test. Such an input receives a complete ordinary double
query, not a nonorientability verdict.

\subsection{Native planning, replay and work accounting}

After the base weighted query completes, discovery reconstructs the finite
graph from the actual base pairings. A parity search either supplies vertex
bits or finds an odd cycle. A zero-parity witness is independently checked
before use. Success constructs the doubled histogram by doubling each
retained multiplicity and skips construction and weighted discovery of the
scale-two system. Failure continues the previous query schedule. The
geometry, boundary selector and base weighted observer are unchanged.

There are $O(t)$ nonempty edge blocks and pairings. The graph implementation
locates the block of each endpoint by binary search, and the parity producer
sorts its vertex keys, giving $O(t\log(t+1))$ indexed arithmetic operations
besides source preparation. Endpoint comparisons have the source's binary
bit bound. The colouring has $O(t\log(t+1))$ encoded bits, and constructing
the histogram takes one checked binary multiplication per retained row.
Neither normal points nor components are expanded. On a miss this finite
work is extra overhead; it is not silently free in the timings. The entire
observer retains its polynomial encoded-input bound. Fewer local queries
do not prove a bound on the outer recognition search or hierarchy states.

The option \code{coorientation=True} is now the default of
\code{normal\_topology\_spectrum}; it requires a strict Boolean. Setting it
false reproduces the old schedule, statistics and certificate exactly.
The CLI exposes \code{--no-coorientation}. Pure vertex-link reduced inputs
still require no orbit query and retain their original certificate.

Successful derivation uses outer schema \code{normal-topology-spectrum-v2},
with the same eight fields as version one. Its query still contains the
boundary transversal, base surface proof, double proof and core spectrum.
The new inner \code{normal-topology-double-coorientation-v1} has exactly
four fields: schema, dimension two, zero-parity colouring and histogram.
It is not an ordinary weighted-orbit certificate and is accepted only through
the source-bound outer protocol. The colouring cannot be transplanted to
another source graph merely because the claimed histogram is plausible.

Independent replay first rebuilds source geometry and canonical reduction,
checks the boundary selector and fully verifies the base weighted proof.
It then reconstructs the graph from that actual base system, checks every
colour constraint, and compares the complete proposed double histogram
with twice the already certified base histogram. It subsequently checks the
spectrum and multiplicity restoration through their existing forward
identities. Replay invokes neither the colouring search, orbit discovery,
weighted producer, sparse inversion nor canonical-core producer. Finite
geometry helpers are shared trusted definitions. Version-one proofs retain
their full explicit double replay; a schema swap cannot turn one proof
language into the other.

Statistics count actual discoveries. Derivation uses two calls rather than
three and records zero double orbit cycles. A shared cycle allowance can
therefore admit a complete derived proof where the old schedule is
inconclusive. Preparation, colouring and verification remain governed by
the cooperative callback, rather than being counted as orbit cycles.
An incomplete base query cannot acquire a derived certificate, and a missed
colouring cannot bypass an exhausted double allowance.

\subsection{Validation and performance}

The accompanying validation compares the complete maintained observer with
a frozen pre-change producer \emph{and} checker, under both periodic rules
and both coordinate-reduction modes. It also compares each complete spectrum
with the independent component records retained in report 53. The disabled
option must reproduce every complete old result, and each failed colouring
must do the same. Successful cases must save exactly the old double cycle
count and return only orientable components. Retained source-bound proofs
are replayed with producer entry points disabled.

Tests independently expand small lifted point systems to check
$w_2(2j+s)=w(j)$ and the generator conjugacy, including reversal bits. They
also cover relabelings, huge binary scalings, orientable coarse misses,
one-sided inputs, strict proof fields, Boolean/integer confusion, altered
weights and multiplicities, duplicate rows, colour flips and missing vertices,
legacy replay, exact cycle allowances and caller cancellation. A simultaneous
colour reversal is accepted because it is a legitimate alternate labelling
of the two sheets.

All 1,322 maintained test methods pass in 223.502 seconds; the focused
39-method run passes in 1.582 seconds. The audit covers 1,275 supplied
vectors under two periodic rules and two reduction modes, or 5,100
configurations. Every full spectrum and source summary agrees with the
frozen \code{ea366dc9b} producer and checker and with the independent
component oracle. All 5,100 disabled-option results reproduce the baseline
exactly. Derivation succeeds on 1,088 configurations and saves precisely
the old double-query cycle count; the other 4,012 return the complete old
result. These counts include scaled and relabelled repeats, not 1,088
distinct knot types. The audit retains 80 paired current/legacy examples
and 557 source hashes.

A separate fresh Regina audit rechecks 664 bounded vectors across 26
triangulation isomorphism classes. All 1,395 produced certificates replay
and match the freshly computed component types. The controls include
orientable genus nine, ten crosscaps and fifteen boundary circles on one
component. Runtime sources remain unchanged during each audit.

The isolated five-round benchmark has 300 measured calls and 60 warmups,
all complete, with interleaved old/old and new/new arms. Timings include
production, independent replay and serialization/hash of the full proof.
The native reference is the frozen producer \emph{and} checker. Each pair
returns the same spectrum hash, and proof bytes are deterministic within
each arm. The following ratios are paired medians.

\begin{samepage}
\noindent\textbf{Full native topology observer.}
\begin{center}
\small
\begin{tabular}{lrrrrr}
\toprule
Input & old ms & new ms & old/new & old A/A & new A/A\\
\midrule
Layered 8 & 9.259 & 4.958 & 1.868 & 1.000 & 0.994\\
Layered 32 & 68.196 & 32.955 & 2.047 & 0.988 & 0.980\\
Layered 64 & 251.061 & 114.663 & 2.190 & 1.017 & 1.002\\
Layered 32, binary scale & 69.058 & 34.356 & 2.004 & 0.997 & 1.003\\
Empty & 0.305 & 0.307 & 0.993 & 0.952 & 0.990\\
Vertex link & 0.325 & 0.324 & 1.007 & 1.012 & 0.970\\
M\"obius band & 1.107 & 1.145 & 0.967 & 0.991 & 1.009\\
Even M\"obius multiple & 1.126 & 1.132 & 0.985 & 0.999 & 0.987\\
Mixed components & 2.218 & 2.228 & 0.995 & 0.994 & 0.978\\
Binary mixture & 2.641 & 3.496 & 0.861 & 0.775 & 0.997\\
Klein bottle & 1.635 & 1.651 & 0.990 & 1.004 & 1.008\\
Layered 256 & 3867.661 & 1704.070 & 2.261 & 0.998 & 0.999\\
\bottomrule
\end{tabular}
\end{center}

\end{samepage}

The layered-256 call improves from about 3.868 to 1.704 seconds, a paired
2.26-fold gain. Its certificate drops from 14,211,485 to 7,139,727 bytes.
The smaller layered gains are 1.87--2.19-fold; the binary-scale case gains
2.00-fold. Pure links and empty reduced inputs keep the identical old
proof and are near parity.

\begin{samepage}
\noindent\textbf{Earlier coordinate-census comparison.}
\begin{center}
\small
\begin{tabular}{lrrrrr}
\toprule
Input & coord. ms & new ms & coord./new & coord. A/A & new A/A\\
\midrule
Layered 8 & 8.789 & 5.036 & 1.758 & 0.983 & 0.963\\
Layered 32 & 62.932 & 32.668 & 1.935 & 0.996 & 0.998\\
Layered 64 & 238.759 & 115.228 & 2.102 & 1.020 & 0.993\\
\bottomrule
\end{tabular}
\end{center}

\end{samepage}

The coordinate reference obtains full component normal vectors, then checks
their topology with the existing scalar observer; it computes finer
embedding data than requested here. This extra work is stated explicitly.
The new topology observer is 1.76--2.10-fold faster on these three primitive
layered controls, resolving the 5--9 percent regressions recorded for the
previous three-query integration. These are supplied-vector calls, not
whole knot-recognition measurements.

The original binary-mixture ratio 0.861 comes with a poor old/old control
of 0.775. The original records are retained. A justified isolated repeat
uses fifteen rounds on that case and two other fallback controls, or 180
measured calls and 12 warmups, all complete. Paired old/new medians are
0.985 for the M\"obius band, 0.993 for the small mixture and 0.994 for
the binary mixture. Repeat A/A medians range from 0.993 to 1.006. The
repeat indicates small extra finite-test costs rather than establishing
a fourteen-percent binary-input regression. It does not erase the initial
noisy measurement or assert zero miss overhead.

\begin{sloppypar}
Raw evidence is in \code{data/weighted-coorientation-*}. From \code{fast/},
run \code{python -B -m} with module
\code{normal\_orbit\_research.}\allowbreak\code{coorientation\_spectra}
with mode \code{audit} or \code{benchmark --rounds 5} and a required
\code{--output} path. The fresh oracle command is
\code{python -B -m} with module \code{topology\_research.regina\_audit}
and mode \code{audit}, with the
retained \code{topology\_research/data/regina\_corpus.json}, an output
path and \code{--recheck-oracle}, in the Regina environment.
\code{data/weighted\_coorientation\_fallback.py} reproduces the longer
isolated control. The table and publication-review scripts retain
completion, source stability, independent replay and document checks.
\end{sloppypar}


This change improves supplied-vector observation and certificate replay.
It adds no diagram recognition coverage and no general quasipolynomial
theorem. Its relevance to the wider programme is that a proof already used
for scalar counts now transports the full additive interface without a
redundant expensive orbit pass. Controlled complete witness discovery and
geometric continuation remain the outstanding global obligations.
